\documentclass[10pt,a4paper]{amsart}
\usepackage[margin=19mm]{geometry}
\usepackage{amsmath,amssymb,mathtools}
\usepackage[scr=rsfs,cal=euler]{mathalfa}
\usepackage{xfrac}
\usepackage[unicode,hidelinks,bookmarks=false]{hyperref}
\newcommand{\ie}{i.e.\ }
\newcommand{\cf}{cf.\ }
\newcommand{\resp}{respectively\ }
\newcommand{\Subsection}{\subsection}
\providecommand{\nobreakdash}{\nobreak\hbox{-}}
\setlength{\emergencystretch}{2em}
\multlinegap0pt
\usepackage{bm}
\usepackage{bbm}
\usepackage{dsfont}
\usepackage{xspace}
\usepackage{cancel}
\usepackage{mathtools}
\usepackage{amsthm}
\makeatletter
\@ifundefined{proposition}{\newtheorem{proposition}{Proposition}}{}
\@ifundefined{lemma}{\newtheorem{lemma}[proposition]{Lemma}}{}
\makeatother
\def\Fbar{\overline{F}}
\def\GM{\mathbb G}
\def\G{G}
\def\bsl{\backslash}
\def\coker{\mathrm{Coker}}
\def\ptf{~.}
\def\tC{{\widetilde C}}
\def\tG{{\widetilde G}}
\def\tI{{\widetilde I}}
\def\tZ{{\widetilde Z}}
\title[Germ expansion for SL(2) in arbitrary characteristics]{Germ expansion for SL(2) in arbitrary characteristics}
\author{Jean-Pierre Labesse}
\date{Source: arXiv:2507.05003v2 (CC BY 4.0)}
\begin{document}
\maketitle

\noindent\emph{Source and licence.} Germ expansion for SL(2) in arbitrary characteristics. Version arXiv:2507.05003v2 (CC BY 4.0), distributed under the Creative Commons Attribution 4.0 International licence, as recorded in the arXiv licence metadata for this exact version and shown on the abstract page https://arxiv.org/abs/2507.05003v2. Full rights notice: licenses/arxiv-2507.05003-CC-BY-4.0.txt.\par\smallskip

\noindent\textit{Editorial scope.} Counted unit: the Lemma whose source label is \texttt{tg} in the article (source0.tex lines 352--355, source label \texttt{tg}), delivered together with the article's own complete original proof of it (source0.tex lines 357--371, one \texttt{proof} environment of 15 lines). No mathematics is written, repaired or re-derived: statement and proof are the article's own text line for line. Required context retained uncounted: none. Every definition, display and label of those lines is retained unchanged. Omitted from this package and not redistributed: 1--351, 356--356, 372--681. Nothing inside a retained excerpt is omitted or altered: every retained body is delivered exactly as the article's lines are written, so no omission receipt of any kind accompanies this package. Citations: the retained excerpts carry no citation command at all, so no bibliography entry of the article is transcribed and no reference is redistributed. Reference adaptation: none was needed and none was made; every reference inside the delivered unit resolves inside it, so no unresolved reference or citation is produced. Printed slips kept literal: no printed word, symbol, hypothesis or number of the counted statement or of its proof was altered. Layout and class: the article's own class is \texttt{amsart} with packages including amssymb, amsmath, amscd, fontenc, enumitem, url, xspace, mathrsfs, euscript, color, xcolor, bbm, bm; the wrapper is the lane's pinned \texttt{amsart} editorial wrapper at 10pt on A4, which restates the theorem-like environments the retained text needs and adds the article's own macro definitions that the retained text needs (\texttt{\textbackslash Fbar}, \texttt{\textbackslash G}, \texttt{\textbackslash GM}, \texttt{\textbackslash bsl}, \texttt{\textbackslash coker}, \texttt{\textbackslash ptf}, \texttt{\textbackslash tC}, \texttt{\textbackslash tG}, \texttt{\textbackslash tI}, \texttt{\textbackslash tZ}), with extra wrapper packages (bm, bbm, dsfont, xspace, cancel, mathtools). The delivered numbering is the unit's own. Assets: no retained span contains a figure environment, a graphics-inclusion command, a PDF-page inclusion, a table or a TikZ picture; no image file is copied into this package, the package declares no asset dependency and no image file is redistributed with it. Licence: version arXiv:2507.05003v2 of this article is distributed under the Creative Commons Attribution 4.0 International licence, as recorded in the arXiv licence metadata for this exact version and shown on the abstract page https://arxiv.org/abs/2507.05003v2. A full rights notice is delivered at licenses/arxiv-2507.05003-CC-BY-4.0.txt. Characters: every retained excerpt is pure ASCII, so the delivered engine, XeLaTeX with its default Latin Modern text font, reports no missing character.
\begin{lemma} \label{tg}
In $\G(F)$ (resp. $\tG(F)$) two elements are stably conjugate if and only if they are
conjugate under $\tG(F)$.
\end{lemma}

\begin{proof}
Firstly one observes that since $\G(\Fbar).\tZ(\Fbar)=\tG(\Fbar)$
conjugacy under $\G(\Fbar)$ and $\tG(\Fbar)$ coincide.
Then it remains to show that conjugacy and stable conjugacy coincide
in $\tG(F)$.  Observe that $\tI_x$ 
is smooth and connected for all $x\in\tG(F)$. In fact 
$\tI_x=Res_{E/F}\GM_m$ for some separable quadratic extension
if $x$ is regular semisimple, $\tI_x=\tG$ if $x$ is central while
 $\tI_x=\GM_m\times\GM_a$ if the Jordan decomposition of $x$ has a non trivial unipotent component.
The set $\tC(x)$ of conjugacy orbits under $\tG(F)$ inside the set
of  rational points in the geometric orbit is parametrized by
 $$\tC(x)\simeq\coker[H^0_{f}(F,\tG)\to H^0_{f}(F,I_x\bsl\tG)]
\simeq\ker[H^1_{f}(F,\tI_x)\to H^1_{f}(F,\tG)]\ptf$$
It remain to observe, that  $H^1_{f}(F,\tI_x)$ is trivial for any $x\in\tG(F)$. 
\end{proof}

\end{document}
